\documentclass[10pt]{article}
\usepackage[T1]{fontenc}
\usepackage[utf8]{inputenc}
\usepackage[french]{babel}
\usepackage{amsmath,amssymb,mathrsfs,mathpazo}
\usepackage{geometry}
\geometry{a4paper,margin=2cm}
\usepackage{graphicx,array,tabularx,tabularray,colortbl,multirow,multicol,booktabs,enumitem,pifont,ulem,fancybox,diagbox,cancel}
\usepackage{pstricks,pstricks-add,pst-plot,pst-tree,pst-node,pst-text,pst-func,pst-math,pst-bspline,pst-3dplot,pst-eucl,pst-solides3d}
\let\pscalculate\undefined
\usepackage{pst-bezier}
\usepackage{tikz,pgfplots,tkz-tab}
\usetikzlibrary{arrows,arrows.meta,patterns,calc,positioning,trees,decorations.markings,shapes.geometric,shapes.misc,shapes.arrows,intersections,decorations.pathmorphing,decorations.pathreplacing,backgrounds,shadows,babel}
\pgfplotsset{compat=1.18}
\usepgfplotslibrary{fillbetween}
\usepackage[np]{numprint}
\usepackage{listings,makecell,pas-tableur,twoopt}
\let\barre\undefined
\usepackage[pstricks]{tablvar}
\definecolor{darkgreen}{RGB}{0,100,0}
\providecommand{\red}{\color{red}}
\providecommand{\blue}{\color{blue}}
\providecommand{\green}{\color{green}}
\providecommand{\black}{\color{black}}
\providecommand{\white}{\color{white}}
\providecommand{\gray}{\color{gray}}
\providecommand{\cyan}{\color{cyan}}
\providecommand{\magenta}{\color{magenta}}
\providecommand{\yellow}{\color{yellow}}
\providecommand{\violet}{\color{violet}}
\providecommand{\orange}{\color{orange}}
\providecommand{\euro}{\texteuro}
\providecommand{\EUR}{\texteuro}
\providecommand{\ieme}{\textsuperscript{e}}
\providecommand{\degres}{\ensuremath{^\circ}}
\providecommand{\degree}{\ensuremath{^\circ}}
\providecommand{\vv}[1]{\overrightarrow{#1}}
\providecommand{\Coord}[1]{#1}
\providecommand{\decosix}{}
\providecommand{\starredbullet}{\ensuremath{\bullet}}
\pagestyle{empty}
\setlength{\parindent}{0pt}
\usepackage{adjustbox}\DeclareUnicodeCharacter{2236}{:}\DeclareUnicodeCharacter{2190}{\ensuremath{\leftarrow}}
\providecommand{\R}{\mathbb{R}}
\providecommand{\N}{\mathbb{N}}
\providecommand{\Z}{\mathbb{Z}}
\providecommand{\Q}{\mathbb{Q}}
\providecommand{\D}{\mathbb{D}}
\providecommand{\C}{\mathbb{C}}
\providecommand{\e}{{\mathrm{e}}}
\providecommand{\vect}[1]{\overrightarrow{#1}}
\providecommand{\vectt}[1]{\overrightarrow{\mathrm{#1}}}
\providecommand{\vvt}[1]{\overrightarrow{\mathrm{#1}}}
\providecommand{\np}[2][]{{#2~#1}}
\providecommand{\vv}[1]{{\overrightarrow{#1}}}
\renewcommand{\d}{{\mathrm{d}}}
\providecommand{\degres}{{^\circ}}
\providecommand{\degre}{{^\circ}}
\providecommand{\blue}{\color{blue}}
\providecommand{\red}{\color{red}}
\providecommand{\enskip}{\quad}
\providecommand{\llbracket}{\lbrack\!\lbrack}
\providecommand{\rrbracket}{\rbrack\!\rbrack}
\providecommand{\euro}{{\text{€}}}
\providecommand{\no}{{\text{n°}}}
\providecommand{\interval}[2]{{[#1;#2]}}
\providecommand{\intervFO}[2]{{[#1;#2[}}
\providecommand{\dotfill}{\ldots}

\providecommand{\ds}{\displaystyle}
\providecommand{\textding}[1]{#1}
\providecommand{\Oij}{\(\mathrm{O};\vec{\imath},\vec{\jmath}\)}
\providecommand{\Oijk}{\(\mathrm{O};\vec{\imath},\vec{\jmath},\vec{k}\)}
\providecommand{\Ouv}{\(\mathrm{O};\vec{u},\vec{v}\)}
\providecommand{\intervFF}[2]{\ensuremath{[#1;#2]}}
\renewcommand{\euro}{\ifmmode\text{\texteuro}\else\texteuro\fi}
\DeclareUnicodeCharacter{00B0}{\ifmmode^{\circ}\else\textdegree\fi}
\DeclareUnicodeCharacter{20AC}{\ifmmode\text{\texteuro}\else\texteuro\fi}
% Énoncés : Éducation nationale. Transcription et corrigés : APMEP, Denis Vergès et contributeurs. Découpage et classement : CoopMaths, Sébastien Lozano, Éric Elter et contributeurs. Adaptation de la présentation : Maths Explorer.
% CC BY-SA : https://coopmaths.fr/www/about/
% https://coopmaths.fr/alea/static/bac/2022/tex/bac_2022_10_sujet1_caledonie_3_cor.tex
\begin{document}
\noindent\textit{APMEP / CoopMaths. CC BY-SA. Présentation adaptée par Maths Explorer.}\par\medskip
\medskip









\begin{center}
\psset{unit=1cm,arrowsize=2pt 3}
\begin{pspicture}(12.2,6)
\pspolygon(0.4,3.4)(0.4,1.5)(1.9,0.4)(1.9,2.3)
\psline(1.9,0.4)(7.9,0.8)(7.9,2.7)(1.9,2.3)
\psline[linestyle=dashed](0.4,1.5)(6.4,1.9)(6.4,3.8)(0.4,3.4)
\psline[linestyle=dashed](7.9,0.8)(6.4,1.9)
\pspolygon(7.9,2.7)(6.4,3.8)(3.4,5.5)(4.9,4.4)
\psline[linestyle=dashed](4.9,4.4)(4.9,2.5)
\psline(4.9,4.4)(1.9,2.3)
\psline(3.4,5.5)(0.4,3.4)
\psline[linewidth=1.5pt]{->}(1.9,0.4)(3.9,0.533)
\psline[linewidth=1.5pt]{->}(1.9,0.4)(1.15,0.95)
\psline[linewidth=1.5pt]{->}(1.9,0.4)(1.9,2.3)
\psline(11.6,0.3)(7.9,1.2)
\psline[linestyle=dashed](7.9,1.2)(7.15,1.35)
\uput[d](1.9,0.4){A} \uput[dr](7.9,0.8){B} \uput[r](6.4,1.9){C} \uput[l](0.4,1.5){D}
\uput[u](1.9,2.3){E} \uput[r](7.9,2.7){F} \uput[ur](6.4,3.8){G} \uput[ul](0.4,3.4){H}
\uput[l](4.9,4.4){I} \uput[u](3.4,5.5){J} \uput[r](11.6,0.3){R} \uput[dl](1.7,0.6){$\vect{\jmath}$}
\uput[dr](2.9,0.45){$\vect{\imath}$} \uput[r](1.9,1.35){$\vect{k}$}
\uput[d](4.9,2.5){L}
\psline(5.1,2.53)(5.1,2.75)(4.9,2.73)
\psdots(11.6,0.3)
\end{pspicture}
\end{center}







\medskip

\begin{enumerate}
\item 
On a $\vect{\text{AG}} = \vect{\text{AB}} + \vect{\text{AD}} + \vect{\text{AE}} = 3\vect{\imath} + 2\vect{\jmath} + \vect{k}$. Donc G(3~;~2~;~1).
\item 


On sait que $M(x~;~y~;~z) \in (\text{EHI}) \iff 2x + 0y - 3z = d$, avec $d \in \R$.

Ainsi, par exemple  E$(0~;~0~;~1) \in (\text{EHI}) \iff 2\times 0 + 0\times 0 - 3\times 1 = d \iff d = -3$.

Donc $M(x~;~y~;~z) \in (\text{EHI}) \iff 2x  - 3z = -3$
\item 
Puisque (EIF) est isocèle en I le projeté orthogonal de I sur [EF] est le milieu de [EF] ; ces deux points ont donc la même abscisse qui est aussi celle du milieu de [AB] soit $\dfrac32$.

L'ordonnée de I est aussi celle de E soit $0$, enfin 

I$\left(\dfrac32~;~0~;~z\right) \in (\text{EHI}) \iff 2 \times \dfrac32  - 3z = - 3\iff 3z = 3 + 3 \iff z = 2$. Donc I$\left(\dfrac32~;~0~;~2\right)$.
\item 
$\bullet~~$ Avec le produit scalaire : avec E(0~;~0~;~1) et F(3~;~0~;~1), on a $\vect{\text{IE}}\begin{pmatrix}-\dfrac32\\0\\- 1\end{pmatrix}$ et $\vect{\text{IF}}\begin{pmatrix}\dfrac32\\0\\- 1\end{pmatrix}$, donc :

IE$^2 = \dfrac94 + 0 + 1 = \dfrac{13}{4}$, d'où IE $ = \dfrac{\sqrt{13}}{2}$ ;

IF$^2 = \dfrac94 + 0 + 1 = \dfrac{13}{4}$, d'où IF $ = \dfrac{\sqrt{13}}{2}$ ;

Donc $\vect{\text{IE}} \cdot \vect{\text{IF}} = \text{IE} \times \text{IF} \times \cos \left(\vect{\text{IF}}~;~ \right)$, soit :

$- \dfrac94 + 0 + 1 = \dfrac{\sqrt{13}}{2} \times \dfrac{\sqrt{13}}{2} \times \cos \left(\vect{\text{IE}}~;~ \vect{\text{IF}}\right) \iff \dfrac{-5}{4} = \dfrac{13}{4}\times \cos \left(\vect{\text{IE}}~;~ \vect{\text{IF}}\right) $.

Finalement $\cos \left(\vect{\text{IE}}~;~ \vect{\text{IF}}\right) = \dfrac{\dfrac{-5}{4}}{\dfrac{13}{4}} = - \dfrac{5}{13}$. La calculatrice donne $\left(\vect{\text{IE}}~;~ \vect{\text{IF}}\right) \approx 112,6\degres$.

$\bullet~~$ Avec le triangle (IEL) rectangle en L (L projeté orthogonal de I sur (EF) :

On a L$\left(\dfrac32~;~0~;~1\right)$, donc IL$^2 = 0 + 0 + 1^2 = 1$, d'où IL $= 1$ ;

On a vu que EI$ = \dfrac{\sqrt{13}}{2}$, d'où $\cos \left(\widehat{\text{EIL}} \right) = \dfrac{\text{IL}}{\text{IE}} = \dfrac{1}{\dfrac{\sqrt{13}}{2}} = \dfrac{2}{\sqrt{13}}$ : la calculatrice donne $\left(\widehat{\text{EIL}} \right) \approx 56,30\degres$ ; donc 
$\widehat{\text{EIF}} = 2\widehat{\text{EIL}}\approx 2 \times 56,30$, soit finalement $\widehat{\text{EIF}} \approx 112,6\degres$.
\item 




	\begin{enumerate}
		\item 
On sait que $M(x~;~y~;~z) \in \Delta \iff \vect{\text{R}M} = t \vect{u} \iff \left\{\begin{array}{l c l}
x - 6&=&- 3t\\
y + 3&=&4t\\
z + 1&=&1t
\end{array}\right.,\:\:  (t \in \R)\iff \left\{\begin{array}{l c l}
x &=&6 - 3t\\
y &=&- 3 + 4t\\
z &=&- 1 + t
\end{array}\right.,\quad  (t \in \R)$.
		\item 
		



K$(x~;~y~;~z)$ est commun à $\Delta$ et au plan (BFG) si ses coordonnées vérifient les équations paramétriques de $\Delta$ et l'équation cartésienne du plan (BFG) soit si le triplet de réels $(x~;~y~;~z)$ vérifie le système :

$\left\{\begin{array}{l c l}
x &=&6 - 3t\\
y &=&- 3 + 4t\\
z &=&- 1 + t\\
x &=& 3
\end{array}\right. \iff \left\{\begin{array}{l c l}
3 &=&6 - 3t\\
y &=&- 3 + 4t\\
z &=&- 1 + t\\
x &=& 3
\end{array}\right. \iff \left\{\begin{array}{l c l}
3t &=&3\\
y &=&- 3 + 4t\\
z &=&- 1 + t\\
x &=& 3
\end{array}\right.\iff$

$\left\{\begin{array}{l c l}
t &=&1\\
y &=&- 3 + 4\\
z &=&- 1 + 1\\
x &=& 3
\end{array}\right.\iff  \left\{\begin{array}{l c l}
t &=&1\\
y &=&1\\
z &=&0\\
x &=&3
\end{array}\right.$

Conclusion : K$(3~;~1~;~0)$.
		\item 
		Avec B(3~;~0~;~0) et C(3~;~2~;~0) on remarque que les coordonnées de K sont les demi-sommes des coordonnées de B et de C, donc que K est le milieu du segment [BC].
	\end{enumerate}
\end{enumerate}

\bigskip

\end{document}
